\documentclass[11pt]{article}
\usepackage[a4paper,margin=1in]{geometry}
\usepackage{amsmath,amssymb,mathtools,amsthm}
\usepackage{enumitem}
\usepackage[bookmarks=false,hidelinks]{hyperref}
\usepackage{microtype}
\allowdisplaybreaks

\newcommand{\F}{\mathbb{F}}
\newcommand{\Z}{\mathbb{Z}}
\newcommand{\suppc}{\operatorname{supp}(\mathbf{c})}
\newcommand{\mymod}{\pmod{m}}
\newcommand{\boldc}{\boldsymbol{c}}
\newcommand{\Moore}{\mathcal{M}oore}
\newtheorem{prop}{Proposition}

\title{Supplementary Material A for\\
\emph{On the Minimum Distance of High-rate Moore-Type Quasi-Cyclic LDPC Codes}\\[0.6ex]
\large Complete Proof of Proposition~29: Classification of Weight-$8$ Matching Configurations}
\author{}
\date{}

\begin{document}
\maketitle

\section{Scope and notation}
This supplementary document contains the complete enumeration proof of Proposition~29 in the main manuscript. It replaces Appendix~IX-A of the submitted version. The notation and case labels are kept identical to those in the main manuscript.

Let $m$ be an odd prime, let $a,b\in\Z_m^*$ satisfy
$\operatorname{ord}_m(a)=k$ and $\operatorname{ord}_m(b)=j$, and assume $j\ge 3$. For the code $\Moore(m,a,b,3)$, a column is represented by
\[
 c(x,t)=\bigl[x,\ x+(b-1)a^t,\ x+(b^2-1)a^t\bigr]^T,
 \qquad x\in\Z_m,\quad t\in\Z_k.
\]
For two support columns $c_i=c(x_i,t_i)$ and $c_q=c(x_q,t_q)$ and a level $s\in\{0,1,2\}$, we write $c_i\sim_s c_q$ if
\[
 x_i+(b^s-1)a^{t_i}\equiv x_q+(b^s-1)a^{t_q}\pmod m.
\]
For each level $s$, the equivalence classes of $\sim_s$ form the matching partition $M_s$, and
$M=(M_0,M_1,M_2)$ is the matching configuration.

The proof below uses the following facts established in the main manuscript.
\begin{enumerate}[label=\textnormal{(\roman*)}]
\item \textbf{Column uniqueness (Lemma~13).} Two represented columns satisfy
$c(x_1,t_1)=c(x_2,t_2)$ if and only if $x_1=x_2$ and $t_1=t_2$.
\item \textbf{Repeated-pair criterion (Lemma~14).} If two distinct support columns match at two distinct levels, then they are equal; hence such a candidate support is impossible.
\item \textbf{Isomorphism reduction (Lemma~19).} Weight-preserving isomorphic matching configurations describe the same support configurations up to a relabeling of the support columns.
\item \textbf{Weight-$8$ block-size restriction (Lemma~20).} At each of the three matching levels, the block-size type is either $2+2+2+2$ or $4+2+2$.
\item \textbf{Normalization (Corollary~23).} Up to an affine permutation preserving the codeword-support property and Hamming weight, one may assume $c_1=c(0,0)$.
\item \textbf{Completeness of the enumeration (Algorithm~1).} Every even matching block admits a refinement into disjoint $2$-matchings. Enumerating all pairwise refinements, all same-level mergings, closing under forced matchings, discarding repeated-column branches, and then reducing by weight-preserving isomorphism yields all legal matching configurations.
\end{enumerate}

\section{Statement of Proposition~29}
\begin{prop}[Proposition~29 of the main manuscript]
Let $\boldc$ be a codeword of $\Moore(m,a,b,3)$ with Hamming weight $8$. Then, without loss of generality, the matching configuration of $\boldc$ is one of the following seven forms.

\begin{enumerate}[label=\textbf{Case \Roman*:},leftmargin=*]
\item
\[
\begin{cases}
M_0=\{\{c_1,c_2\},\{c_3,c_4\},\{c_5,c_6\},\{c_7,c_8\}\},\\
M_1=\{\{c_1,c_3\},\{c_2,c_4\},\{c_5,c_7\},\{c_6,c_8\}\},\\
M_2=\{\{c_1,c_4,c_6,c_7\},\{c_2,c_8\},\{c_3,c_5\}\}.
\end{cases}
\]
\item
\[
\begin{cases}
M_0=\{\{c_1,c_2\},\{c_3,c_4\},\{c_5,c_6\},\{c_7,c_8\}\},\\
M_1=\{\{c_1,c_3,c_6,c_8\},\{c_2,c_5\},\{c_4,c_7\}\},\\
M_2=\{\{c_1,c_4\},\{c_2,c_3\},\{c_5,c_8\},\{c_6,c_7\}\}.
\end{cases}
\]
\item
\[
\begin{cases}
M_0=\{\{c_1,c_2\},\{c_3,c_4,c_5,c_6\},\{c_7,c_8\}\},\\
M_1=\{\{c_1,c_3\},\{c_2,c_5\},\{c_4,c_7\},\{c_6,c_8\}\},\\
M_2=\{\{c_1,c_4\},\{c_2,c_6\},\{c_3,c_7\},\{c_5,c_8\}\}.
\end{cases}
\]
\item
\[
\begin{cases}
M_0=\{\{c_1,c_2\},\{c_3,c_4\},\{c_5,c_6\},\{c_7,c_8\}\},\\
M_1=\{\{c_1,c_3\},\{c_2,c_4\},\{c_5,c_7\},\{c_6,c_8\}\},\\
M_2=\{\{c_1,c_5\},\{c_2,c_7\},\{c_3,c_6\},\{c_4,c_8\}\}.
\end{cases}
\]
\item
\[
\begin{cases}
M_0=\{\{c_1,c_2\},\{c_3,c_4\},\{c_5,c_6\},\{c_7,c_8\}\},\\
M_1=\{\{c_1,c_3\},\{c_2,c_4\},\{c_5,c_7\},\{c_6,c_8\}\},\\
M_2=\{\{c_1,c_5\},\{c_2,c_6\},\{c_3,c_7\},\{c_4,c_8\}\}.
\end{cases}
\]
\item
\[
\begin{cases}
M_0=\{\{c_1,c_2\},\{c_3,c_4\},\{c_5,c_6\},\{c_7,c_8\}\},\\
M_1=\{\{c_1,c_3\},\{c_2,c_6\},\{c_5,c_7\},\{c_4,c_8\}\},\\
M_2=\{\{c_1,c_5\},\{c_2,c_4\},\{c_3,c_7\},\{c_6,c_8\}\}.
\end{cases}
\]
\item
\[
\begin{cases}
M_0=\{\{c_1,c_2\},\{c_3,c_4\},\{c_5,c_6\},\{c_7,c_8\}\},\\
M_1=\{\{c_1,c_3\},\{c_2,c_7\},\{c_4,c_5\},\{c_6,c_8\}\},\\
M_2=\{\{c_1,c_5\},\{c_2,c_6\},\{c_3,c_7\},\{c_4,c_8\}\}.
\end{cases}
\]
\end{enumerate}
\end{prop}

\section{Complete enumeration proof}
\begin{proof}
Let $\boldc$ be a codeword of $\Moore(m,a,b,3)$ of Hamming weight $8$. By Corollary~23 of the main manuscript, we may assume without loss of generality that $c_1=c(0,0)$.

By Lemma~20 of the main manuscript, at each matching level only the block-size types $2+2+2+2$ and $4+2+2$ can occur. By Algorithm~1 of the main manuscript, it is sufficient to begin with pairwise refinements; the possible and forced mergings are treated below.

\emph{Enumeration procedure.} Every even matching block admits a pairwise refinement. We enumerate these refinements, include same-level mergings, and close each candidate under the matchings forced by its congruences, as in Algorithm~1 of the main manuscript. At each stage, repeated-column tests discard impossible branches, and relabelings preserving the fixed levels identify equivalent choices. The calculations below give the surviving configurations; the final isomorphism reduction identifies seven representatives.

After relabeling the support columns, we may refine the matching at level $s=0$ as
$$
M_0=
\bigl\{
\{c_1,c_2\},
\{c_3,c_4\},
\{c_5,c_6\},
\{c_7,c_8\}
\bigr\}.
$$
Equivalently, the first row of the support matrix is written as
$$
0,0,x,x,y,y,z,z.
$$
This is only a pairwise refinement of the matching at level $s=0$;
we do not assume that $0,x,y,z$ are pairwise distinct. Thus the
support matrix has the form
$$
\suppc=[c_1,c_2,\ldots,c_8],
$$
where
$$
\begin{aligned}
    \suppc
    =&\left[
    \begin{array}{cccc}
        0 & 0 & x & x \\
        b-1 & (b-1)a^{t_2}
        & x+(b-1)a^{t_3} & x+(b-1)a^{t_4} \\
        b^2-1 & (b^2-1)a^{t_2}
        & x+(b^2-1)a^{t_3} & x+(b^2-1)a^{t_4}
    \end{array}
    \right. \\
    &\left.
    \begin{array}{cccc}
        y & y & z & z \\
        y+(b-1)a^{t_5} & y+(b-1)a^{t_6}
        & z+(b-1)a^{t_7} & z+(b-1)a^{t_8} \\
        y+(b^2-1)a^{t_5} & y+(b^2-1)a^{t_6}
        & z+(b^2-1)a^{t_7} & z+(b^2-1)a^{t_8}
    \end{array}
    \right].
\end{aligned}
$$

By Lemma~14 of the main manuscript, if the same pair of distinct support columns is matched at two different levels, then the two columns coincide. Therefore, every branch containing such a repeated pair is discarded.

At each step below, we enumerate all admissible pairwise matchings which contain no pair already matched at a previous level. The residual relabeling symmetries preserving the previously fixed matchings are used to identify equivalent choices.

Since the labels of the support columns are arbitrary, we may assume
that $c_1$ is matched with $c_3$ at level $s=1$. At level $s=2$, the
column $c_1$ cannot be matched with $c_2$ or $c_3$. Up to the
remaining relabeling symmetries, there are only two possibilities:
either $c_1$ is matched with the other column in the same $M_0$-block
as $c_3$, or it is matched with a column in another $M_0$-block.
Hence we obtain the following two principal cases:
\begin{align*}
    \text{Case 1}
    \begin{cases}
    	b-1 &\equiv x+(b-1)a^{t_3}\mymod, \\
    	b^2-1 &\equiv x+(b^2-1)a^{t_4}\mymod,
    \end{cases}
    \text{Case 2}
    \begin{cases}
    	b-1 &\equiv x+(b-1)a^{t_3}\mymod, \\
    	b^2-1 &\equiv y+(b^2-1)a^{t_5}\mymod.
    \end{cases}
    \end{align*}
We now discuss the possible matchings at levels $s=1$ and $s=2$ in these two cases.
\paragraph{Case 1}
In this case, the support matrix becomes
\begin{eqnarray*}
		~&~&\left[\begin{array}{cccc}
			0 & 0 
			& x & x \\
			b-1 & (b-1)a^{t_2} 
			& b-1 & (b-1)+\frac{b}{b+1}x\\
            b^2 -1 & (b^2-1)a^{t_2} 
			& b^2-1-bx &b^2-1
            \end{array}
            \right.\\
			~&~& \left. \begin{array}{cccc}
             y & y 
			& z & z\\
			  y+(b-1)a^{t_5} & y+(b-1)a^{t_6}
			& z+(b-1)a^{t_7} & z+(b-1)a^{t_8} \\
			  y+(b^2-1)a^{t_5} & y+(b^2-1)a^{t_6}
			& z+(b^2-1)a^{t_7} & z+(b^2-1)a^{t_8}
        \end{array}
        \right].
\end{eqnarray*}
Since $j\geq 3$, we have $b+1\not\equiv0\mymod$, so the
above expressions are well defined.

At level $s=1$, up to relabeling, there are two possible matchings which contain no pair already matched at level $s=0$:
\[
\text{Case 1-1:}\qquad
\begin{cases}
(b-1)+\dfrac{b}{b+1}x \equiv (b-1)a^{t_2}\mymod,\\
y+(b-1)a^{t_5} \equiv z+(b-1)a^{t_7}\mymod,\\
y+(b-1)a^{t_6} \equiv z+(b-1)a^{t_8}\mymod.
\end{cases}
\]
\[
\text{Case 1-2:}\qquad
\begin{cases}
(b-1)+\dfrac{b}{b+1}x \equiv z+(b-1)a^{t_7}\mymod,\\
(b-1)a^{t_2} \equiv y+(b-1)a^{t_5}\mymod,\\
y+(b-1)a^{t_6} \equiv z+(b-1)a^{t_8}\mymod.
\end{cases}
\]
\paragraph{Case 1-1}
The support matrix becomes
\begin{eqnarray*}
		~&~&\left[\begin{array}{ccccc}
			0 & 0 & x & x & y\\
			  b-1 & b-1+\frac{b}{b+1}x & b-1 & b-1+\frac{b}{b+1}x & y+(b-1)a^{t_5}\\
            b^2 -1 & b^2-1+bx & b^2-1-bx &b^2-1 & y+(b^2-1)a^{t_5}
            \end{array} \right. \\
		~&~& \left. \begin{array}{ccc} 
            y & z & z\\
		      y+(b-1)a^{t_6} & y+(b-1)a^{t_5} & y+(b-1)a^{t_6} \\
			y+(b^2-1)a^{t_6} & (b+1)y-bz+(b^2-1)a^{t_5} & (b+1)y-bz+(b^2-1)a^{t_6}
        \end{array}
        \right]
        .
\end{eqnarray*}
At level $s=2$, if
$b^2-1+bx\equiv b^2-1-bx$,
then $x\equiv0\mymod$, which forces $c_1=c_3$, a contradiction. 
Hence, up to relabeling, the only remaining subcase is the following.

\paragraph{Case 1-1-1}
$$
\begin{cases}
	b^2-1-bx &\equiv y+(b^2-1)a^{t_5}\mymod,\\
	(b+1)y-bz+(b^2-1)a^{t_5} &\equiv y+(b^2-1)a^{t_6}\mymod,\\
	(b+1)y-bz+(b^2-1)a^{t_6} &\equiv b^2-1+bx\mymod.
\end{cases}
$$
Adding the three congruences gives $y\equiv x+z \mymod$.
Substituting this into the third congruence gives
$y+(b^2-1)a^{t_6}\equiv b^2-1\mymod$, i.e., $c_6\sim_2 c_1$.
Thus a matching block of size four occurs at level $s=2$, namely
$$
\{c_1,c_4,c_6,c_7\}\in M_2.
$$
This gives a legal matching configuration.

\paragraph{Case 1-2}
The support matrix becomes
\begin{eqnarray*}
		~&~&\left[
        \begin{array}{ccccc}
			0 & 0 & x & x & y\\
            b-1 & (b-1)a^{t_2} & b-1 & b-1+\frac{b}{b+1}x & (b-1)a^{t_2}\\
            b^2 -1 &  (b^2-1)a^{t_2} & b^2-1-bx & b^2-1 & -by+(b^2-1)a^{t_2}
        \end{array}
        \right.\\
		~&~& \left. 
        \begin{array}{cccc}
			y & z & z\\			
			  y+(b-1)a^{t_6} & b-1+\frac{b}{b+1}x  & y+(b-1)a^{t_6} \\			
			y+(b^2-1)a^{t_6} & bx-bz+b^2-1 & (b+1)y-bz+(b^2-1)a^{t_6}
        \end{array}
        \right]
        .
\end{eqnarray*}
At level $s=2$, up to relabeling, there are six possible subcases.

\paragraph{Case 1-2-1}
$$
\begin{cases}
	b^2-1-bx &\equiv  (b^2-1)a^{t_2}\mymod,\\
	-by+(b^2-1)a^{t_2} &\equiv (b+1)y-bz+(b^2-1)a^{t_6}\mymod, \\
	y+(b^2-1)a^{t_6} &\equiv bx-bz+b^2-1\mymod.
\end{cases}
$$
Adding the three congruences gives
$2b(z-y-x)\equiv0\mymod$,
and hence $z\equiv x+y\mymod$. 
Substituting this into the third congruence gives 
$y+(b-1)a^{t_6}\equiv b-1 \mymod$, i.e., $c_6\sim_1 c_1$.
Thus a matching block of size four occurs at level $s=1$:
$$
\{c_1,c_3,c_6,c_8\}\in M_1.
$$
This gives a legal matching configuration.

\paragraph{Case 1-2-2}
$$
\begin{cases}
	y+(b^2-1)a^{t_6}&\equiv (b^2-1)a^{t_2}\mymod, \\
	b^2-1-bx &\equiv bx-bz+b^2-1\mymod, \\
	-by+(b^2-1)a^{t_2} &\equiv (b+1)y-bz+(b^2-1)a^{t_6}\mymod.
\end{cases}
$$
The second congruence gives $z\equiv2x\mymod$, and adding the first and the
third congruences gives $z\equiv2y\mymod$. Hence $x\equiv y\mymod$.
Therefore a matching block of size four occurs at level $s=0$:
$$
\{c_3,c_4,c_5,c_6\}\in M_0.
$$
This gives a legal matching configuration.

\paragraph{Case 1-2-3}
$$
\begin{cases}
	y+(b^2-1)a^{t_6} &\equiv (b^2-1)a^{t_2}\mymod,\\
	b^2-1-bx &\equiv (b+1)y-bz+(b^2-1)a^{t_6}\mymod,\\
	-by+(b^2-1)a^{t_2} &\equiv bx-bz+b^2-1\mymod.
\end{cases}
$$
Adding the three congruences gives $z\equiv x+y\mymod$. Substituting this
into the third congruence gives $a^{t_2}\equiv1\mymod$, and hence
$c_1=c_2$ by Lemma~13 of the main manuscript, a contradiction.

\paragraph{Case 1-2-4}
$$
\begin{cases}
	(b^2-1)a^{t_2} &\equiv bx-bz+b^2-1\mymod, \\
	b^2-1-bx &\equiv y+(b^2-1)a^{t_6}\mymod,\\
    (b+1)y-bz+(b^2-1)a^{t_6}&\equiv -by+(b^2-1)a^{t_2}\mymod.
\end{cases}
$$
Adding the three congruences gives $x\equiv y\mymod$. Substituting this
into the second congruence gives $y+(b-1)a^{t_6}\equiv b-1\mymod$, i.e. $c_6\sim_1 c_1$.
Since $c_1 \sim_1 c_3$ and $c_3\sim_2 c_6$,
we obtain $c_3=c_6$ from Lemma~14 of the main manuscript, a contradiction.

\paragraph{Case 1-2-5}
$$
\begin{cases}
	 (b+1)y-bz+(b^2-1)a^{t_6}&\equiv (b^2-1)a^{t_2}\mymod, \\
	-by+(b^2-1)a^{t_2} &\equiv b^2-1-bx\mymod,\\
	bx-bz+b^2-1 &\equiv y+(b^2-1)a^{t_6}\mymod.
\end{cases}
$$
Adding the three congruences gives $x\equiv z\mymod$, i.e., $\{c_3,c_4,c_7,c_8\}\in M_0$, which forces $c_4=c_7$ from $c_4\sim_0 c_7$ and $c_4\sim_1 c_7$ by Lemma~14 of the main manuscript, a contradiction.

\paragraph{Case 1-2-6}
$$
\begin{cases}
	(b+1)y-bz+(b^2-1)a^{t_6} &\equiv (b^2-1)a^{t_2}\mymod,  \\
	-by+(b^2-1)a^{t_2} &\equiv bx-bz+b^2-1\mymod, \\
	b^2-1-bx&\equiv y+(b^2-1)a^{t_6}\mymod.
\end{cases}
$$
Adding the three congruences gives $x\equiv0\mymod$, i.e., $\{c_1,c_2,c_3,c_4\}\in M_0$ which forces
$c_1=c_3$ from $c_1\sim_0 c_3$ and $c_1\sim_1 c_3$ by Lemma~14 of the main manuscript, a contradiction.

\paragraph{Case 2}
In this case, the support matrix becomes
\begin{eqnarray*}
		~&~&\left[\begin{array}{cccc}
			0 & 0 
			& x & x \\
            b-1 & (b-1)a^{t_2} 
			& b-1 & x+(b-1)a^{t_4}\\
            b^2 -1 & (b^2-1)a^{t_2} 
			& b^2-1-bx & x+(b^2-1)a^{t_4}
            \end{array}
            \right.\\
			~&~& \left. \begin{array}{cccc}            
			y & y 
			& z & z\\			
			  b-1+\frac{b}{b+1}y  & y+(b-1)a^{t_6}
			& z+(b-1)a^{t_7} & z+(b-1)a^{t_8} \\
			  b^2-1 & y+(b^2-1)a^{t_6}
			& z+(b^2-1)a^{t_7} & z+(b^2-1)a^{t_8}
        \end{array}
        \right]
        .
\end{eqnarray*}
At level $s=1$, up to relabeling, there are five possible matchings which contain no pair already matched at level $s=0$:
\[
\text{Case 2-1:}\qquad
\begin{cases}
(b-1)a^{t_2} \equiv x+(b-1)a^{t_4}\mymod,\\
b-1+\dfrac{b}{b+1}y \equiv z+(b-1)a^{t_7}\mymod,\\
y+(b-1)a^{t_6} \equiv z+(b-1)a^{t_8}\mymod.
\end{cases}
\]
\[
\text{Case 2-2:}\qquad
\begin{cases}
(b-1)a^{t_2} \equiv b-1+\dfrac{b}{b+1}y\mymod,\\
x+(b-1)a^{t_4} \equiv z+(b-1)a^{t_7}\mymod,\\
y+(b-1)a^{t_6} \equiv z+(b-1)a^{t_8}\mymod.
\end{cases}
\]
\[
\text{Case 2-3:}\qquad
\begin{cases}
(b-1)a^{t_2} \equiv y+(b-1)a^{t_6}\mymod,\\
x+(b-1)a^{t_4} \equiv z+(b-1)a^{t_7}\mymod,\\
b-1+\dfrac{b}{b+1}y \equiv z+(b-1)a^{t_8}\mymod.
\end{cases}
\]
\[
\text{Case 2-4:}\qquad
\begin{cases}
(b-1)a^{t_2} \equiv z+(b-1)a^{t_7}\mymod,\\
x+(b-1)a^{t_4} \equiv b-1+\dfrac{b}{b+1}y\mymod,\\
y+(b-1)a^{t_6} \equiv z+(b-1)a^{t_8}\mymod.
\end{cases}
\]
\[
\text{Case 2-5:}\qquad
\begin{cases}
(b-1)a^{t_2} \equiv z+(b-1)a^{t_7}\mymod,\\
z+(b-1)a^{t_8} \equiv b-1+\dfrac{b}{b+1}y\mymod,\\
x+(b-1)a^{t_4} \equiv y+(b-1)a^{t_6}\mymod.
\end{cases}
\]

\paragraph{Case 2-1}
The support matrix becomes
\begin{eqnarray*}
		~&~&\left[\begin{array}{cccc}
			0 & 0 
			& x & x \\
            b-1 & (b-1)a^{t_2} 
			& b-1 & (b-1)a^{t_2} \\
            b^2 -1 & (b^2-1)a^{t_2} 
			& b^2-1-bx & -bx+(b^2-1)a^{t_2}
            \end{array}
            \right.\\
			~&~& \left. \begin{array}{cccc}   
			y & y 
			& z & z\\			
			b-1+\frac{b}{b+1}y  & y+(b-1)a^{t_6}
			& b-1+\frac{b}{b+1}y & y+(b-1)a^{t_6}\\
			b^2-1 & y+(b^2-1)a^{t_6}
			& by-bz+b^2-1 & (b+1)y-bz+(b^2-1)a^{t_6}
        \end{array}
        \right]
        .
\end{eqnarray*}
At level $s=2$, up to relabeling, there are seven possible subcases.

\paragraph{Case 2-1-1}
$$
\begin{cases}
	b^2-1-bx &\equiv (b^2-1)a^{t_2}\mymod, \\
	-bx+(b^2-1)a^{t_2} &\equiv (b+1)y-bz+(b^2-1)a^{t_6}\mymod, \\
	y+(b^2-1)a^{t_6} &\equiv by-bz+b^2-1\mymod.
\end{cases}
$$
Adding the three congruences gives $z\equiv x+y\mymod$. Substituting this into the first and third congruences gives
$y+(b^2-1)a^{t_6}\equiv (b^2-1)a^{t_2}\mymod$, i.e., $c_6\sim_2 c_2$. Therefore a matching block of size four occurs at level $s=2$:
$$
\{c_2,c_3,c_6,c_7\}\in M_2.
$$
This gives a legal matching configuration.

\paragraph{Case 2-1-2}
$$
\begin{cases}
	b^2-1-bx &\equiv y+(b^2-1)a^{t_6}\mymod, \\
	(b^2-1)a^{t_2} &\equiv by-bz+b^2-1\mymod, \\
	(b+1)y-bz+(b^2-1)a^{t_6} &\equiv -bx+(b^2-1)a^{t_2}\mymod.
\end{cases}
$$
This case does not force repeated support columns, and hence it gives a legal matching configuration.

\paragraph{Case 2-1-3}
$$
\begin{cases}
	b^2-1-bx &\equiv y+(b^2-1)a^{t_6}\mymod,\\
	(b+1)y-bz+(b^2-1)a^{t_6}&\equiv  (b^2-1)a^{t_2}\mymod,\\
	-bx+(b^2-1)a^{t_2} &\equiv by-bz+b^2-1\mymod.
\end{cases}
$$
Adding the three congruences gives $-2bx\equiv0\mymod$, and hence
$x\equiv0\mymod$. Thus $c_1=c_3$ from $c_1\sim_0 c_3$ and $c_1\sim_1 c_3$ by Lemma~14 of the main manuscript, a contradiction.

\paragraph{Case 2-1-4}
$$
\begin{cases}
	b^2-1-bx &\equiv by-bz+b^2-1\mymod, \\
	  y+(b^2-1)a^{t_6}&\equiv -bx+(b^2-1)a^{t_2}\mymod, \\
	(b^2-1)a^{t_2} &\equiv (b+1)y-bz+(b^2-1)a^{t_6}\mymod.
\end{cases}
$$
Adding the three congruences gives $0\equiv2b(y-z)\mymod$, and hence
$z\equiv y\mymod$. Thus $c_6=c_8$ from $c_6\sim_0 c_8$ and $c_6\sim_1 c_8$ by Lemma~14 of the main manuscript, a contradiction.

\paragraph{Case 2-1-5}
$$
\begin{cases}
	b^2-1-bx &\equiv (b+1)y-bz+(b^2-1)a^{t_6}\mymod, \\
	(b^2-1)a^{t_2} &\equiv by-bz+b^2-1\mymod, \\
	y+(b^2-1)a^{t_6} &\equiv -bx+(b^2-1)a^{t_2}\mymod.
\end{cases}
$$
Adding the three congruences gives $0\equiv2b(y-z)\mymod$, and hence
$z\equiv y$. Similarly, $c_6=c_8$, a contradiction.

\paragraph{Case 2-1-6}
$$
\begin{cases}
	b^2-1-bx &\equiv (b+1)y-bz+(b^2-1)a^{t_6}\mymod, \\
	y+(b^2-1)a^{t_6} &\equiv (b^2-1)a^{t_2}\mymod,\\
	-bx+(b^2-1)a^{t_2} &\equiv by-bz+b^2-1\mymod.
\end{cases}
$$
Adding the three congruences gives $-2bx\equiv2b(y-z)\mymod$, and hence $z\equiv x+y\mymod$. Substituting this into the first congruence gives 
$y+(b^2-1)a^{t_6}\equiv b^2-1$, i.e., $c_6\sim_2 c_1$.
Since $c_1\sim_2 c_5$ and $c_5\sim_0 c_6$, we obtain $c_5=c_6$ by Lemma~14 of the main manuscript, a contradiction.

\paragraph{Case 2-1-7}
$$
\begin{cases}
	-bx+(b^2-1)a^{t_2} &\equiv (b+1)y-bz+(b^2-1)a^{t_6}\mymod,\\
	y+(b^2-1)a^{t_6} &\equiv (b^2-1)a^{t_2}\mymod, \\
	b^2-1-bx &\equiv by-bz+b^2-1\mymod.
\end{cases}
$$
Adding the three congruences gives $-2bx\equiv2b(y-z)\mymod$, and hence $z\equiv x+y\mymod$. This case does not force repeated support columns, and therefore it gives a legal matching  configuration.

\paragraph{Case 2-2}
The support matrix becomes
\begin{eqnarray*}
		~&~&\left[\begin{array}{ccccc}
			0 & 0 & x & x & y\\
            b-1 & b-1+\frac{b}{b+1}y & b-1 & x+(b-1)a^{t_4} & b-1+\frac{b}{b+1}y \\
            b^2 -1 & by+b^2-1 & b^2-1-bx & x+(b^2-1)a^{t_4} & b^2-1
            \end{array}
            \right.\\
			~&~& \left. \begin{array}{ccc}              
			y & z & z\\			
			  y+(b-1)a^{t_6} & x+(b-1)a^{t_4}  & y+(b-1)a^{t_6}\\	
            y+(b^2-1)a^{t_6} & (b+1)x-bz+(b^2-1)a^{t_4} & (b+1)y-bz+(b^2-1)a^{t_6}
        \end{array}
        \right]
        .
\end{eqnarray*}
At level $s=2$, up to relabeling, there are six possible subcases.

\paragraph{Case 2-2-1}
$$
\begin{cases}
	b^2-1-bx &\equiv by+b^2-1\mymod, \\
	x+(b^2-1)a^{t_4} &\equiv (b+1)y-bz+(b^2-1)a^{t_6}\mymod,\\
	y+(b^2-1)a^{t_6} &\equiv (b+1)x-bz+(b^2-1)a^{t_4}\mymod.
\end{cases}
$$
The first congruence gives $x+y\equiv0\mymod$, while adding the three congruences gives $z\equiv x+y\mymod$. Hence $z\equiv0\mymod$. Therefore a matching block of size four occurs at level $s=0$:
$$
\{c_1,c_2,c_7,c_8\}\in M_0.
$$
This gives a legal matching configuration.

\paragraph{Case 2-2-2}
$$
\begin{cases}
	b^2-1-bx &\equiv y+(b^2-1)a^{t_6}\mymod, \\
	(b+1)y-bz+(b^2-1)a^{t_6} &\equiv x+(b^2-1)a^{t_4}\mymod, \\
	(b+1)x-bz+(b^2-1)a^{t_4} &\equiv by+b^2-1\mymod.
\end{cases}
$$
Adding the three congruences gives $-2bz\equiv0$, and hence $z\equiv0\mymod$. Then the pair $\{c_2,c_7\}$ is matched at $s=0,2$. By Lemma~14 of the main manuscript, this forces $c_2=c_7$, a
contradiction.

\paragraph{Case 2-2-3}
$$
\begin{cases}
	b^2-1-bx &\equiv (b+1)x-bz+(b^2-1)a^{t_4}\mymod, \\
	y+(b^2-1)a^{t_6} &\equiv by+(b^2-1)\mymod, \\
	x+(b^2-1)a^{t_4} &\equiv (b+1)y-bz+(b^2-1)a^{t_6}\mymod.
\end{cases}
$$
Adding the three congruences gives $0\equiv2b(x+y-z)\mymod$, and hence $z\equiv x+y\mymod$. Substituting this into the first congruence and dividing by $b+1$ gives
$$
b-1+\dfrac{b}{b+1}y \equiv x+(b-1)a^{t_4}\mymod.
$$
Therefore a matching block of size four occurs at level $s=1$:
$$
\{c_2,c_4,c_5,c_7\}\in M_1.
$$
This gives a legal matching configuration.

\paragraph{Case 2-2-4}
$$
\begin{cases}
	b^2-1-bx &\equiv (b+1)x-bz+(b^2-1)a^{t_4}\mymod, \\
	 x+(b^2-1)a^{t_4} &\equiv y+(b^2-1)a^{t_6}\mymod,\\
	(b+1)y-bz+(b^2-1)a^{t_6} &\equiv by+(b^2-1)\mymod.
\end{cases}
$$
Adding the three congruences gives $-2bx\equiv0\mymod$, and hence
$x\equiv0\mymod$. Then the pair $\{c_1,c_3\}$ is matched at levels $s=0,1$. By Lemma~14 of the main manuscript, this forces $c_1=c_3$, a
contradiction.

\paragraph{Case 2-2-5}
$$
\begin{cases}
	b^2-1-bx &\equiv (b+1)y-bz+(b^2-1)a^{t_6}\mymod, \\
	x+(b^2-1)a^{t_4} &\equiv by+(b^2-1)\mymod, \\
	y+(b^2-1)a^{t_6} &\equiv (b+1)x-bz+(b^2-1)a^{t_4}\mymod.
\end{cases}
$$
Adding the three congruences gives $0\equiv2b(x+y-z)\mymod$, and hence $z\equiv x+y\mymod$. 
Substituting this into the first congruence gives 
$$
b^2-1\equiv y+(b^2-1)a^{t_6}\mymod.
$$
Then the pair $\{c_5,c_6\}$ is matched at levels $s=0,2$. By
Lemma~14 of the main manuscript, this forces $c_5=c_6$, a contradiction.

\paragraph{Case 2-2-6}
$$
\begin{cases}
	b^2-1-bx &\equiv (b+1)y-bz+(b^2-1)a^{t_6}\mymod, \\
	y+(b^2-1)a^{t_6} &\equiv  x+(b^2-1)a^{t_4}\mymod,\\
	(b+1)x-bz+(b^2-1)a^{t_4} &\equiv by+b^2-1\mymod.
\end{cases}
$$
Adding the three congruences gives $2by\equiv0$, and hence $y\equiv0\mymod$. 
Then the pair $\{c_2,c_5\}$ is matched at levels $s=0,2$. By Lemma~14 of the main manuscript, this forces $c_2=c_5$, a
contradiction.

\paragraph{Case 2-3}
The support matrix becomes:
\begin{eqnarray*}
		~&~&\left[\begin{array}{ccccc}
			0 & 0 & x & x & y\\
			b-1 & (b-1)a^{t_2} & b-1 & x+(b-1)a^{t_4} & b-1+\frac{b}{b+1}y\\
            b^2 -1 & (b^2-1)a^{t_2} & b^2-1-bx & x+(b^2-1)a^{t_4} & b^2-1 
            \end{array}
            \right.\\
		~&~& \left. \begin{array}{ccc}    
			y & z & z\\ 
			  (b-1)a^{t_2} & x+(b-1)a^{t_4}  & b-1+\frac{b}{b+1}y\\
			-by+(b^2-1)a^{t_2} & (b+1)x-bz+(b^2-1)a^{t_4} & by-bz+b^2-1
        \end{array}
        \right]
        .
\end{eqnarray*}
At level $s=2$, up to relabeling, there are six possible subcases.

\paragraph{Case 2-3-1}
$$
\begin{cases}
	b^2-1-bx &\equiv (b^2-1)a^{t_2}\mymod, \\
	x+(b^2-1)a^{t_4} &\equiv by-bz+b^2-1\mymod, \\
	-by+(b^2-1)a^{t_2} &\equiv (b+1)x-bz+(b^2-1)a^{t_4}\mymod.
\end{cases}
$$
Adding the three congruences gives $0\equiv2b(x+y-z)\mymod$, and hence $z\equiv x+y\mymod$. 
Substituting this into the first congruence gives
$$
b^2-1-bx\equiv by-bz+b^2-1\equiv x+(b^2-1)a^{t_4}\mymod.
$$
Then the pair $\{c_3,c_4\}$ is matched at levels $s=0,2$. By
Lemma~14 of the main manuscript, this forces $c_3=c_4$, a contradiction.

\paragraph{Case 2-3-2}
$$
\begin{cases}
	b^2-1-bx &\equiv -by+(b^2-1)a^{t_2}\mymod, \\
	x+(b^2-1)a^{t_4} &\equiv by-bz+b^2-1\mymod, \\
	(b^2-1)a^{t_2}&\equiv (b+1)x-bz+(b^2-1)a^{t_4}\mymod.
\end{cases}
$$
Adding the three congruences gives $0\equiv2b(x-z)\mymod$, and hence
$x\equiv z\mymod$. 
Then the pair $\{c_4,c_7\}$ is matched at levels $s=0,1$. By Lemma~14 of the main manuscript, this forces $c_4=c_7$, a contradiction.

\paragraph{Case 2-3-3}
$$
\begin{cases}
	b^2-1-bx &\equiv (b+1)x-bz+(b^2-1)a^{t_4}\mymod,\\
	 x+(b^2-1)a^{t_4}&\equiv (b^2-1)a^{t_2}\mymod,\\
	-by+(b^2-1)a^{t_2} &\equiv by-bz+b^2-1\mymod.
\end{cases}
$$
Adding the three congruences gives $0\equiv2b(x+y-z)\mymod$, and hence $z\equiv x+y\mymod$. Moreover, the first congruence gives
$$
(b+1)x-bz+(b^2-1)a^{t_4}
\equiv b^2-1-bx
\equiv by-bz+b^2-1\mymod,
$$
i.e. $c_7\sim_2 c_8$.
Because $c_7\sim_0 c_8$ and $c_7\sim_2 c_8$, this forces $c_7=c_8$ by Lemma~14 of the main manuscript, a contradiction.

\paragraph{Case 2-3-4}
$$
\begin{cases}
	b^2-1-bx &\equiv (b+1)x-bz+(b^2-1)a^{t_4}\mymod, \\
	(b^2-1)a^{t_2} &\equiv by-bz+b^2-1\mymod, \\
	  x+(b^2-1)a^{t_4}&\equiv -by+(b^2-1)a^{t_2}\mymod.
\end{cases}
$$
Adding the three congruences gives $0\equiv2b(x-z)\mymod$, and hence
$x\equiv z\mymod$. Then the pair 
$\{c_4,c_7\}$ is matched at levels $s=0,1$. By Lemma~14 of the main manuscript, this forces $c_4=c_7$, a
contradiction.

\paragraph{Case 2-3-5}
$$
\begin{cases}
	b^2-1-bx &\equiv by-bz+b^2-1\mymod, \\
	(b^2-1)a^{t_2} &\equiv x+(b^2-1)a^{t_4}\mymod, \\
	(b+1)x-bz+(b^2-1)a^{t_4} &\equiv -by+(b^2-1)a^{t_2}\mymod.
\end{cases}
$$
The first congruence gives $z\equiv x+y\mymod$. This case does not force repeated support columns, and therefore it gives a legal matching configuration.

\paragraph{Case 2-3-6}
$$
\begin{cases}
	b^2-1-bx &\equiv by-bz+b^2-1\mymod,\\
	(b^2-1)a^{t_2} &\equiv (b+1)x-bz+(b^2-1)a^{t_4}\mymod, \\
	   x+(b^2-1)a^{t_4}&\equiv -by+(b^2-1)a^{t_2}\mymod.
\end{cases}
$$
The first congruence gives $z\equiv x+y\mymod$. Adding the three congruences gives $0\equiv2b(x-z)\mymod$, and hence $x\equiv z\mymod$. Therefore $y\equiv0\mymod$, i.e., $c_1\sim_0 c_5$. Then the pair $\{c_1,c_5\}$ is matched at levels $s=0,2$. By Lemma~14 of the main manuscript, this forces $c_1=c_5$, a contradiction.

\paragraph{Case 2-4}
The support matrix becomes
\begin{eqnarray*}
		~&~&\left[\begin{array}{ccccc}
        	0 & 0 & x & x & y\\
            b-1 & (b-1)a^{t_2} & b-1 & b-1+\frac{b}{b+1}y & b-1+\frac{b}{b+1}y \\
            b^2 -1 & (b^2-1)a^{t_2} & b^2-1-bx & by-bx+b^2-1 & b^2-1 
            \end{array}
            \right.\\
            ~&~& \left. \begin{array}{ccc}            
        	y & z & z\\        	
        	y+(b-1)a^{t_6} & (b-1)a^{t_2}   & y+(b-1)a^{t_6}\\ 
        	y+(b^2-1)a^{t_6} & -bz+(b^2-1)a^{t_2} & (b+1)y-bz+(b^2-1)a^{t_6}
        \end{array}
        \right]
        .
\end{eqnarray*}
At level $s=2$, up to relabeling, there are six possible subcases.

\paragraph{Case 2-4-1}
$$
\begin{cases}
	b^2-1-bx &\equiv (b^2-1)a^{t_2}\mymod, \\
	-bz+(b^2-1)a^{t_2} &\equiv y+(b^2-1)a^{t_6}\mymod, \\
	 (b+1)y-bz+(b^2-1)a^{t_6}&\equiv by-bx+b^2-1\mymod.
\end{cases}
$$
Adding the three congruences gives $-2bz\equiv0\mymod$, and hence
$z\equiv0\mymod$, i.e., $c_2\sim_0 c_7$. Then the pair $\{c_2,c_7\}$ is matched at two distinct levels $s=0,1$. By Lemma~14 of the main manuscript, this forces $c_2=c_7$, a
contradiction.

\paragraph{Case 2-4-2}
$$
\begin{cases}
	 y+(b^2-1)a^{t_6}&\equiv b^2-1-bx\mymod, \\
	(b^2-1)a^{t_2} &\equiv (b+1)y-bz+(b^2-1)a^{t_6} \mymod,\\
	by-bx+b^2-1 &\equiv -bz+(b^2-1)a^{t_2}\mymod.
\end{cases}
$$
Adding the three congruences gives $0\equiv-2bz\mymod$, and hence
$z\equiv0\mymod$. Similarly, this forces $c_2=c_7$, a
contradiction.

\paragraph{Case 2-4-3}
$$
\begin{cases}
	  b^2-1-bx &\equiv-bz+(b^2-1)a^{t_2}\mymod,\\
	(b^2-1)a^{t_2} &\equiv y+(b^2-1)a^{t_6}\mymod, \\
	(b+1)y-bz+(b^2-1)a^{t_6} &\equiv by-bx+b^2-1\mymod.
\end{cases}
$$
This case does not force repeated support columns, and therefore it
gives a legal matching configuration.

\paragraph{Case 2-4-4}
$$
\begin{cases}
	b^2-1-bx &\equiv -bz+(b^2-1)a^{t_2}\mymod, \\
	y+(b^2-1)a^{t_6} &\equiv by-bx+b^2-1\mymod,  \\
	(b^2-1)a^{t_2} &\equiv (b+1)y-bz+(b^2-1)a^{t_6}\mymod.
\end{cases}
$$
Adding the three congruences gives $0\equiv2b(y-z)\mymod$, and hence
$y\equiv z\mymod$, i.e., $c_6\sim_0 c_8$. Then the pair $\{c_6,c_8\}$ is matched at $s=0,1$. By Lemma~14 of the main manuscript, this forces $c_6=c_8$, a contradiction.

\paragraph{Case 2-4-5}
$$
\begin{cases}
	b^2-1-bx &\equiv (b+1)y-bz+(b^2-1)a^{t_6}\mymod, \\
	(b^2-1)a^{t_2} &\equiv by-bx+b^2-1\mymod, \\
	y+(b^2-1)a^{t_6} &\equiv -bz+(b^2-1)a^{t_2}\mymod.
\end{cases}
$$
Adding the three congruences gives $0\equiv2b(y-z)\mymod$, and hence
$y\equiv z\mymod$. Similarly, by Lemma~14 of the main manuscript, this forces $c_6=c_8$, a
contradiction.

\paragraph{Case 2-4-6}
$$
\begin{cases}
	b^2-1-bx &\equiv (b+1)y-bz+(b^2-1)a^{t_6}\mymod, \\
	y+(b^2-1)a^{t_6} &\equiv (b^2-1)a^{t_2}\mymod, \\
	-bz+(b^2-1)a^{t_2} &\equiv by-bx+b^2-1\mymod.
\end{cases}
$$
Adding the three congruences gives $0\equiv2by\mymod$, and hence
$y\equiv0\mymod$. Then the pair $\{c_1,c_5\}$ is matched at levels $s=0,2$. By Lemma~14 of the main manuscript, this forces $c_1=c_5$, a
contradiction.

\paragraph{Case 2-5}
The support matrix becomes
\begin{eqnarray*}
		~&~&\left[\begin{array}{cccc}
			0 & 0 &
			x & x \\
            b-1 & (b-1)a^{t_2} &
			b-1 & x+(b-1)a^{t_4} \\
            b^2-1 & (b^2-1)a^{t_2} &
			b^2-1-bx & x+(b^2-1)a^{t_4} \\
            \end{array}
            \right.\\
            ~&~& \left. \begin{array}{cccc}  
			y & y &
			z & z \\			
			b-1+\frac{b}{b+1}y & x+(b-1)a^{t_4} &
			(b-1)a^{t_2} & b-1+\frac{b}{b+1}y \\			
			b^2-1 & (b+1)x-by+(b^2-1)a^{t_4} &
			-bz+(b^2-1)a^{t_2} & by-bz+b^2-1
        \end{array}
        \right]
        .
\end{eqnarray*}
At level $s=2$, up to relabeling, there are seven possible subcases.

\paragraph{Case 2-5-1}
$$
\begin{cases}
b^2-1-bx &\equiv (b^2-1)a^{t_2}\mymod, \\
-bz+(b^2-1)a^{t_2} &\equiv x+(b^2-1)a^{t_4}\mymod, \\
(b+1)x-by+(b^2-1)a^{t_4} &\equiv by-bz+b^2-1\mymod.
\end{cases}
$$
Adding the three congruences gives $0\equiv2by\mymod$, and hence
$y\equiv0\mymod$. Then the pair $\{c_1,c_5\}$ is matched at levels $s=0,2$. By Lemma~14 of the main manuscript, this forces $c_1=c_5$, a
contradiction.

\paragraph{Case 2-5-2}
$$
\begin{cases}
    b^2-1-bx &\equiv (b^2-1)a^{t_2}\mymod,\\
    x+(b^2-1)a^{t_4} &\equiv by-bz+b^2-1\mymod,\\
	-bz+(b^2-1)a^{t_2} &\equiv (b+1)x-by+(b^2-1)a^{t_4}\mymod.
\end{cases}
$$
Adding the three congruences gives $0\equiv2bx\mymod$, and hence
$x\equiv0\mymod$. Then the pair $\{c_1,c_3\}$ is matched at levels $s=0,1$. By Lemma~14 of the main manuscript, this forces $c_1=c_3$, a
contradiction.

\paragraph{Case 2-5-3}
$$
\begin{cases}
	b^2-1-bx &\equiv  (b+1)x-by+(b^2-1)a^{t_4}\mymod,\\
	(b^2-1)a^{t_2} &\equiv by-bz+b^2-1\mymod, \\
	x+(b^2-1)a^{t_4}&\equiv -bz+(b^2-1)a^{t_2}\mymod.
\end{cases}
$$
Adding the three congruences gives $0\equiv2b(x-z)\mymod$, and hence
$x\equiv z\mymod$. Then the pair $\{c_4,c_7\}$ is matched at levels $s=0,2$. By Lemma~14 of the main manuscript, this forces $c_4=c_7$, a
contradiction.

\paragraph{Case 2-5-4}
$$
\begin{cases}
	b^2-1-bx &\equiv -bz+(b^2-1)a^{t_2}\mymod, \\
	(b^2-1)a^{t_2} &\equiv x+(b^2-1)a^{t_4}\mymod, \\
	(b+1)x -by+(b^2-1)a^{t_4} &\equiv by-bz+b^2-1\mymod.
\end{cases}
$$
Adding the three congruences gives $2by\equiv2bz\mymod$, and hence
$y\equiv z\mymod$. Then the pair $\{c_5,c_8\}$ is matched at levels $s=0,1$. By Lemma~14 of the main manuscript, this forces $c_5=c_8$, a
contradiction.

\paragraph{Case 2-5-5}
$$
\begin{cases}
	b^2-1-bx &\equiv -bz+(b^2-1)a^{t_2}\mymod, \\
	x+(b^2-1)a^{t_4} &\equiv by-bz+b^2-1\mymod, \\
	  (b^2-1)a^{t_2}&\equiv (b+1)x-by+(b^2-1)a^{t_4}\mymod.
\end{cases}
$$
Adding the three congruences gives $0\equiv2b(x-z)\mymod$, and hence
$x\equiv z\mymod$. Substituting this into the first congruence gives
$$
b^2-1\equiv (b^2-1)a^{t_2}\mymod,
$$
i.e., $c_1\sim_2 c_2$.
Then the pair $\{c_1,c_2\}$ is matched at $s=0,2$. By Lemma~14 of the main manuscript, this forces $c_1=c_2$, a contradiction.

\paragraph{Case 2-5-6}
$$
\begin{cases}
	b^2-1-bx &\equiv by-bz+b^2-1\mymod, \\
	(b^2-1)a^{t_2} &\equiv x+(b^2-1)a^{t_4}\mymod, \\
	(b+1)x-by+(b^2-1)a^{t_4} &\equiv -bz+(b^2-1)a^{t_2}\mymod.
\end{cases}
$$
The first congruence gives $z\equiv x+y\mymod$. Adding the three
congruences gives $0\equiv2b(y-z)$, and hence $y\equiv z\mymod$.
Therefore $x\equiv0\mymod$. Then the pair $\{c_1,c_3\}$ is matched at levels $s=0,1$. By Lemma~14 of the main manuscript, this forces $c_1=c_3$, a contradiction.

\paragraph{Case 2-5-7}
$$
\begin{cases}
	b^2-1-bx &\equiv by-bz+b^2-1\mymod, \\
	x+(b^2-1)a^{t_4} &\equiv -bz+(b^2-1)a^{t_2}\mymod, \\
	(b^2-1)a^{t_2} &\equiv (b+1)x-by+(b^2-1)a^{t_4}\mymod.
\end{cases}
$$
The first congruence gives $z\equiv x+y\mymod$. Adding the three
congruences gives $0\equiv2b(x-z)\mymod$, and hence $x\equiv z\mymod$.
Therefore $y\equiv0\mymod$. Then the pair $\{c_1,c_5\}$ is matched at levels $s=0,2$. By Lemma~14 of the main manuscript, this forces
$c_1=c_5$, a contradiction.

The enumeration above contains $39$ nontrivial subcases. Every
discarded subcase forces either repeated support columns or a pair of
distinct support columns matched at two different levels. The
remaining subcases give the following ten legal matching
configurations:
$$
1\text{-}1\text{-}1,\quad
1\text{-}2\text{-}1,\quad
1\text{-}2\text{-}2,\quad
2\text{-}1\text{-}1,\quad
2\text{-}1\text{-}2,\quad
2\text{-}1\text{-}7,\quad
2\text{-}2\text{-}1,\quad
2\text{-}2\text{-}3,\quad
2\text{-}3\text{-}5,\quad
2\text{-}4\text{-}3.
$$

The forced $4$-matchings have been retained throughout the enumeration. Together with the pairwise-refinement procedure, this accounts for both block-size types in Lemma~20 of the main manuscript.

Among the ten matching configurations, consider the permutations
$$
\pi_1=(1\,2)(3\,4)(5\,8)(6\,7),
$$
$$
\pi_2=(1\,2)(3\,5)(4\,6)(7\,8),
$$
$$
\pi_3=(1\,3)(2\,4)(5\,7)(6\,8).
$$
A direct verification shows that $\pi_1$ is a weight-preserving
graph isomorphism between the matching graphs of Cases 1-1-1
and 2-1-1, $\pi_2$ is a weight-preserving graph isomorphism
between those of Cases 1-2-1 and 2-2-3, and $\pi_3$ is a
weight-preserving graph isomorphism between those of Cases
1-2-2 and 2-2-1.

By Lemma~19 of the main manuscript, each of these pairs describes the same support configurations up to a weight-preserving relabeling.

Consequently, it is sufficient to retain the seven representatives
$$
1\text{-}1\text{-}1,\quad
1\text{-}2\text{-}1,\quad
1\text{-}2\text{-}2,\quad
2\text{-}1\text{-}2,\quad
2\text{-}1\text{-}7,\quad
2\text{-}3\text{-}5,\quad
2\text{-}4\text{-}3.
$$
We denote the above seven cases by Case I--Case VII, respectively.
This proves Proposition~29 of the main manuscript.
\end{proof}
\section{Conclusion}
The preceding case analysis exhausts the pairwise refinements and all admissible same-level mergings prescribed by the enumeration procedure. Every discarded branch forces repeated support columns, whereas the surviving branches reduce, under weight-preserving relabeling, to exactly the seven representatives listed in Proposition~29. This completes the classification of weight-$8$ matching configurations.

\end{document}
